\documentclass[12pt]{article}
\usepackage[utf8]{inputenc}
\usepackage[T1]{fontenc}
\usepackage[margin=1in]{geometry}
\usepackage{setspace}
\usepackage{xcolor}
\usepackage[colorlinks=true, urlcolor=blue!60!black, linkcolor=blue!60!black]{hyperref}
\setlength{\parindent}{0em}
\setlength{\parskip}{0.6em}
\renewcommand{\baselinestretch}{1.0}
\begin{document}
\begin{center}
{\LARGE \textit{Calder v. Bull}, 3 U.S. (3 Dall.) 386 (1798)\par}
\vspace{0.5em}
{\large Case Brief --- Two-Page Edition\par}
\end{center}

\section{Memory Jogger}
The Ex Post Facto Clause bars only retroactive criminal laws; Chase and Iredell debated natural law versus written text (\href{https://supreme.justia.com/cases/federal/us/3/386/}{Justia}).

\section{Detailed Case Facts}
A Connecticut probate court disapproved Normand Morrison's will, so his estate passed to the Calders by intestacy. After the Bulls' time to appeal lapsed, the legislature---which by long usage exercised judicial powers---set aside the decree and granted a new hearing, at which the will was approved. The Calders called the resolution an \textit{ex post facto} law.

\section{Procedural History}
Connecticut's Superior Court and Supreme Court of Errors upheld the resolution; the Calders brought a writ of error.

\section{Judicial Votes}
Affirmed 4--0 in \textit{seriatim} opinions by Chase, Paterson, Iredell, and Cushing; Ellsworth and Wilson did not participate.

\section{Holding}
The Ex Post Facto Clauses reach only retroactive criminal laws, so Connecticut's civil resolution was valid. The Court cannot void a state law for violating the state's own constitution. The natural-law debate was dicta.

\section{Analysis of Opinions}
\textbf{Chase:} ``vital principles in our free republican governments'' limit legislatures even without express text, but ``ex post facto'' is a technical term for penal laws; otherwise the Contracts and Takings Clauses would be superfluous. \textbf{Paterson:} the term meant ``crimes, pains, and penalties,'' though he wished it reached all retrospective laws. \textbf{Iredell:} courts may void only laws that violate the written Constitution, never laws thought contrary to natural justice. \textbf{Cushing:} the resolution was valid whether judicial or legislative.

\section{Examples: Future Applications}
\textbf{Same-Side:} a retroactive civil tax amendment is permitted. \textbf{Opposite-Side:} retroactively dropping a corroboration requirement is barred (\textit{Carmell v. Texas}). \textbf{Fence-Sitter:} retroactive sex-offender registration is valid if genuinely civil (\textit{Smith v. Doe}), but punitive residency bans may cross the line.

\section{Critique}
Chase condemned taking property from A to give to B, yet upheld a resolution that shifted an inheritance on a vesting technicality. Originalists praise Iredell's positivism, while Justice Thomas has questioned the criminal-only limitation.

\section{Key Quotations}
\begin{quote}
``1st. Every law that makes an action done before the passing of the law and which was innocent when done, criminal and punishes such action. 2d. Every law that aggravates a crime, or makes it greater than it was when committed. 3d. Every law that changes the punishment, and inflicts a greater punishment than the law annexed to the crime, when committed. 4th. Every law that alters the legal rules of evidence and receives less or different testimony than the law required at the time of the commission of the offence in order to convict the offender.''
\end{quote}
\noindent --- Justice Chase, 3 U.S. at 390.

\end{document}
